% §15.4 Ranked Search Application
%
% Cross-reference resolution status:
%
% DEFINED here:
%   \label{sec:ranked-search}                    -- this section
%   \label{prop:bphf-ranked-search}              -- expected-miss bound
%   \label{eq:bphf-ranked-misses}                -- expected miss count
%
% ALREADY DEFINED (chapter 5):
%   \Cref{def:bloom-filter}                      -- chapter 5 §5.2
%
% ALREADY DEFINED (chapter 6):
%   \Cref{def:perfect-filter-line}               -- chapter 6 §6.4
%
% ALREADY DEFINED (chapter 14):
%   \Cref{def:bhf}                               -- chapter 14 §14.2
%   \Cref{prop:bhf-fpr}                          -- chapter 14 §14.4
%
% ALREADY DEFINED (chapter 15):
%   \Cref{def:bphf}                              -- §15.2
%   \Cref{sec:bphf-construction}                 -- §15.2
%   \Cref{alg:bphf-build-query}                  -- §15.2
%   \Cref{prop:bphf-space}                       -- §15.3
%   \Cref{sec:space-time-analysis}               -- §15.3
%   \eqref{eq:bphf-fnr}                          -- §15.2 FNR bound
%   \eqref{eq:bphf-r}                            -- §15.2 rank width
%
% FORWARD REFS within chapter 15:
%   \Cref{sec:bridge-to-part5}                   -- §15.5 (Task 6)
%
% §15.4 develops the BPHF's archetypal application: top-k retrieval.
% The ranked-search problem is stated; the BPHF is constructed by
% pre-computing scores, sorting, and assigning ranks; the structural
% argument for why FPR = 0 is the right guarantee in this setting is
% given; the expected-miss bound is derived; the worked example at
% n = 10^6 and k = 100 makes the bit budget concrete.

\section{Ranked Search Application}
\label{sec:ranked-search}

The chapter has so far developed the BPHF as a membership filter
on the FPR-zero arm of the perfect-filter line. The structural
features (zero false positives in the witness query model, exact at
full resolution with a true perfect hash, and a controlled false
negative rate $\fnrate \approx 1 - m/n$ bought by a $k$-perfect hash
with $m$ slots, at per-element cost $\lceil \log_2 n \rceil$ bits from
\Cref{prop:bphf-space}) describe the filter in the abstract; the
present section places it in its archetypal application, which is also
the setting that supplies the witness model the FPR-zero guarantee
needs. The setting is ranked search: given a corpus and a
query, return the top-$k$ most relevant documents. The exact
solution scores every document and sorts; the BPHF accelerates the
solution by storing each document's pre-computed rank in a perfect-
hash-addressed table, with the query checking whether a candidate
document's stored rank falls within the top-$k$ band. The
$\fprate = 0$ guarantee is exactly the precision invariant the
application needs, and the chapter argues that the BPHF is not a
generic data structure that happens to fit ranked search but the
data structure ranked search asks for.

\subsection*{The ranked-search problem}

Fix a corpus of $n$ documents indexed by IDs $a_1, \ldots, a_n$. For
each document the application has a relevance scoring function
$s : \text{ID} \to \mathbb{R}$, computed from features such as
term-frequency weights, semantic embeddings, or external metadata.
A query asks for the top-$k$ documents, meaning the $k$ document IDs
whose scores are largest under $s$. The exact solution computes
$s(a_i)$ for every $i$, sorts the scores in descending order, and
returns the IDs at positions $1, \ldots, k$. The cost is $O(n)$
score evaluations plus $O(n \log n)$ sorting, dominated by the score
evaluations when $s$ is expensive. At corpus sizes in the millions,
the per-query cost is dominated by the score-evaluation phase, and
the application's question becomes: can the top-$k$ retrieval avoid
evaluating $s$ on every document at every query? An indexing
structure that pre-computes the answer once and serves queries
without per-document scoring is the right architectural answer; the
BPHF is one such structure.

\subsection*{BPHF for top-\texorpdfstring{$k$}{k} retrieval}

The construction is a single-pass build over the corpus. Pre-compute
$s(a_i)$ for every document, sort the IDs in descending order of
score, and assign to each ID its rank in the sort: the ID with the
largest score gets rank $0$, the next gets rank $1$, and so on, with
the ID at the bottom of the sort getting rank $n - 1$. The rank
assignment $\rho : A \to \{0, 1, \ldots, n - 1\}$ is exactly the
function the BPHF construction of \Cref{sec:bphf-construction}
requires as input. Build the BPHF on the corpus with this $\rho$ and
$r = \lceil \log_2 n \rceil$ bits per slot from \eqref{eq:bphf-r}.
The resulting data structure has the property that for any document
ID $a \in A$, the query $\textsc{Query}(a)$ from
\Cref{alg:bphf-build-query} reads $R[\phi(a)] = \rho(a)$ and returns
$\top$, recovering the document's pre-computed rank.

The top-$k$ retrieval reduces to a rank-bound check. For a candidate
document ID $a$, the application asks: is $a$ in the top $k$? The
top-$k$ predicate is $\rho(a) < k$ (since the top-scoring document
has rank $0$ and the $k$-th has rank $k - 1$). The query algorithm
becomes
\[
  \textsc{TopK}(a) \;=\; \textsc{Query}(a) \;\land\; R[\phi(a)] < k,
\]
combining the BPHF's membership check with a numeric comparison
against $k$. The membership check returns $\top$ when $a$ is a known
corpus document; the numeric comparison admits $a$ only when its
stored rank is in the top band. Both conditions must hold for the
top-$k$ verdict to be $\top$.

The query cost is $O(1)$: one perfect-hash evaluation $\phi(a)$, one
rank-function evaluation $\rho(a)$, one memory read of $R[\phi(a)]$,
one $r$-bit comparison against $\rho(a)$ (the BPHF check), and one
integer comparison against $k$ (the top-$k$ check). The per-query
work is constant in $n$ and constant in $k$; the application can
serve top-$k$ queries at the rate of one per few hundred CPU cycles.
The build cost is the one-time work of scoring all $n$ documents,
sorting them, and constructing the BPHF, totaling $O(n \log n)$
sorting plus the $O(n)$ BPHF build that \Cref{sec:space-time-analysis}
analyzed. The amortization is favorable: each query saves $n$ score
evaluations relative to the exact solution, so the per-query cost
drops from $O(n)$ to $O(1)$ once the build is paid off.

The architecture is suitable for any application that issues many
queries against a stable corpus. The build runs once when the
corpus is loaded or refreshed; the queries run continuously
afterwards. A search engine ingests a new corpus snapshot every
hour and rebuilds the BPHF; in between rebuilds the BPHF serves
millions of queries per second on commodity hardware. A
recommendation system ingests user-rated items nightly and rebuilds
the per-user BPHFs; in between rebuilds the personalized top-$k$
predictions are served instantly. The pattern, batch build and
constant-time query, matches the BHF's deployment profile from
chapter~14 and applies here to the FPR-zero side of the design
space.

\subsection*{Why \texorpdfstring{$\fprate = 0$}{FPR = 0} is the right guarantee}

The top-$k$ retrieval problem has an asymmetric error structure
that distinguishes the two error modes by consequence. A false
positive is the BPHF reporting that an irrelevant or off-topic
document belongs in the top set; a false negative is the BPHF
missing a document that should have been included. The two errors
corrupt the application's output in different ways.

A false positive corrupts the result list directly. The application
shows the user a document that has no business being there; the
user sees an obvious irrelevance among the top results; the
precision of the result list visibly suffers. In a search interface
the user clicks the irrelevant result, finds it useless, and judges
the search system as broken. In a recommendation system the
irrelevant item gets shown, the user dismisses it, and the
recommendation engine's trust score drops. The cost of a false
positive is the cost of corrupted output: visible to the user,
attributable to the system, and immediately damaging to the
application's perceived quality.

A false negative has a softer consequence. The application omits a
document that should have been included; the user sees a result
list with one entry missing; the recall of the result list suffers.
In many settings the user does not notice (a top-$100$ list is
indistinguishable from a top-$99$-plus-one-missing list to the
casual reader), and even when the user does notice, the missing
entry is recovered by relaxing the query, browsing further, or
re-running with looser parameters. The cost of a false negative is
the cost of mild under-recall: invisible to most users, hard to
attribute to the system, and recoverable through application-level
adaptations.

The asymmetry tracks the BPHF's error structure exactly. The
$\fprate = 0$ guarantee from \Cref{def:bphf} eliminates the
high-cost error mode by construction; the $\fnrate > 0$ trade-off
admits the low-cost error mode at a controlled rate. The
application designer who chooses the BPHF over a competing filter
makes the trade explicitly: zero corrupted output in exchange for a
small under-recall. The trade is the right one for any setting where
precision is the binding quality metric and recall is the
recoverable one. Ranked search, content-recommendation systems,
precision-sensitive classification, and authorization filters are
all examples; the section returns to the broader category in the
discussion that follows.

\subsection*{The false-negative cost}

Ranked search is the witness model: a query is posed for a document
that already carries its corpus rank, and a non-corpus candidate has
no such rank to present. The BPHF therefore returns no off-corpus
document, $\fprate = 0$. The false negatives come entirely from the
$k$-perfect hash. At full resolution, $m = n$ slots and a true
perfect hash, every corpus document owns its slot and the retrieval
is exact: no top-$k$ document is missed. Shrinking the table to
$m < n$ slots crowds out a fraction $1 - m/n$ of the corpus, and the
expected number of top-$k$ documents missed is that fraction of the
$k$ retrieved:
\begin{equation}
\label{eq:bphf-ranked-misses}
  \mathbb{E}[\text{misses}]
  \;\approx\; k \cdot \fnrate
  \;\approx\; k \left( 1 - \frac{m}{n} \right),
\end{equation}
using $\fnrate \approx 1 - m/n$ from \eqref{eq:bphf-fnr} and
linearity of expectation over the $k$ true top documents. The miss
count is set by the retrieval depth $k$ and the table-size ratio
$m/n$, the designer's space-versus-recall dial, and is independent of
the rank width $r$, which is fixed at $\lceil \log_2 n \rceil$ by the
need for distinct ranks.

\begin{proposition}[BPHF top-$k$ expected misses]
\label{prop:bphf-ranked-search}
Let the BPHF of \Cref{def:bphf} be built on a corpus of $n$
documents in the witness model, with a $k$-perfect hash of $m \le n$
slots. A top-$k$ retrieval query against this BPHF surfaces no
off-corpus document ($\fprate = 0$) and misses, in expectation, at
most $k\,(1 - m/n)$ of the true top-$k$ documents. At full
resolution $m = n$ the retrieval is exact.
\end{proposition}

The proposition matches the intuition that a fuller table misses
fewer documents and a deeper retrieval exposes more documents to the
miss rate. Halving the spare capacity $1 - m/n$ halves the expected
misses; setting $m = n$ removes them entirely. The designer picks the
table size to make the per-query miss expectation small enough that
the user-facing experience is indistinguishable from exact top-$k$
retrieval, paying in slots for the recall it wants.

\subsection*{A worked example}

Fix a corpus of $n = 10^6$ documents and a top-$k = 100$ retrieval
target. The rank width is
\[
  r \;=\; \lceil \log_2 10^6 \rceil
  \;=\; \lceil 19.93 \rceil
  \;=\; 20 \text{ bits},
\]
since $2^{20} = 1{,}048{,}576 \ge 10^6 > 2^{19} = 524{,}288$. The
BPHF's per-element bit cost from \Cref{prop:bphf-space} is therefore
$20$ bits plus the sublinear perfect-hash overhead. The rank table
occupies
\[
  n \cdot r \;=\; 10^6 \cdot 20 \;=\; 2 \cdot 10^7 \text{ bits}
  \;=\; 2.5 \text{ megabytes}.
\]
Adding the perfect-hash representation (between $2$ and $3$ bits per
element at $n = 10^6$ in practical implementations), the total is
approximately $2.7$ to $2.9$ megabytes. The whole indexing structure
fits comfortably in the L3 cache of a modern server CPU, and the
per-query work is a handful of memory accesses entirely served from
cache.

At full resolution, $m = n = 10^6$ slots with a true perfect hash,
the error guarantees are
\[
  \fprate \;=\; 0
  \qquad \text{and} \qquad
  \fnrate \;=\; 0,
\]
and the retrieval is exact: every query returns the right $100$
documents, never an off-corpus one. To trade table space for the
usual reason, shrink to a $k$-perfect hash with $m = 0.99\,n$ slots.
By \eqref{eq:bphf-ranked-misses} the expected miss count is
\[
  \mathbb{E}[\text{misses}]
  \;\approx\; 100 \left( 1 - \frac{0.99\,n}{n} \right)
  \;=\; 100 \cdot 0.01 \;=\; 1,
\]
so a typical top-$100$ query returns about $99$ of the true top
documents, while
\[
  \fprate \;=\; 0
\]
still holds exactly: not one of the roughly $|\Universe| - n$
off-corpus candidates is ever surfaced. The one-percent table saving
costs about one missed result in a hundred, and the affirmative
guarantee, never showing an irrelevant document, is untouched. From
the user's perspective the retrieval is indistinguishable from exact
top-$k$: the missing entry is rarely noticed, and a wrong entry never
appears.

The bit cost matches the application's needs. A baseline solution
that stores the full sorted index would need $n \log_2 n = 10^6
\cdot 20 = 2 \cdot 10^7$ bits for the ranks plus the perfect-hash
addressing, the same $2.5$ to $3$ megabytes as the BPHF. The
baseline gives exact retrieval at the same space cost, so what does
the BPHF add? The answer is in the query algorithm. The baseline
needs to support efficient lookup of a document's rank by ID, which
requires either a hash table (with collision handling, load factor
overhead, and potentially $O(1)$ amortized but variable-time
queries) or a sorted associative structure (with $O(\log n)$
queries). The BPHF provides the same lookup at one memory access
and one comparison, the minimum any data structure can achieve. The
$0$ false positive guarantee is the bonus: the BPHF refuses
non-corpus queries by construction, where a hash table would need
an additional membership check to filter them out.

The construction also extends to a useful family of generalizations.
A multi-resolution BPHF stores ranks at varying precisions per
document (more bits for high-ranked documents, fewer bits for the
tail), reducing the per-element cost for applications that only
query the head of the distribution. A streaming-rebuild BPHF
maintains the rank index as new documents arrive, by amortizing the
rebuild cost over the inter-arrival interval. A tiered BPHF
combines a primary BPHF on the recent corpus with a secondary BPHF
on the archived corpus, with queries checking both and merging the
results. Each generalization preserves the $\fprate = 0$ guarantee
of the underlying BPHF and varies the operational profile to suit a
specific deployment. The exercises return to these variants; the
present section establishes the baseline construction and its
error structure, and the chapter's bridge in
\Cref{sec:bridge-to-part5} returns to the place of the BPHF in
Part~IV's broader architecture.

The ranked-search application illustrates the larger point of
Part~IV: the perfect-filter line of \Cref{def:perfect-filter-line}
is not a theoretical curiosity but an architectural decision space
that maps cleanly to operational requirements. Applications with
precision-binding quality metrics belong on the FPR-zero arm and
use the BPHF; applications with recall-binding quality metrics
belong on the FNR-zero arm and use the BHF. The decision is forced
by the application's tolerance for each error mode, and the
chapter's two constructions cover both halves of the resulting
design space at the asymptotic bit-budget optimum on each side.
